% Solutions to the problems of Chapter 5 (Autonomous Driving as a Representative Application).
% Included by backmatter/solutions.tex.  Numbers come from the scripts in
% companion_repo/exercises/ch05/solutions/ and companion_repo/examples/ch05_*/.
\begin{addedblock}
\begin{sol}{prob:ch05:1}
(a)~For a point near the optical axis, $\Delta u\approx f_{x}\,\delta\psi=1200\times8.73\times10^{-3}=10.5$ pixels, independent of the range: a rotation of the camera rotates the whole image. (b)~With \eqref{eq:yaw-error-lateral}, $\Delta y\approx r\,\delta\psi$: $0.17$\,m at $20$\,m and $0.70$\,m at $80$\,m, so the metric error grows linearly with range while the pixel error does not. (c)~A translation $\delta t=0.1$\,m along the camera $x$ axis moves the projection by $\Delta u=f_{x}\,\delta t/Z$: $6.0$ pixels at $20$\,m and $1.5$ pixels at $80$\,m, whereas the metric error is $0.10$\,m at both ranges. Translation errors therefore dominate near-field fusion (parking, cut-ins at short gap) and rotation errors dominate far-field fusion (highway detection at $60$--$150$\,m), which is why a camera--LiDAR calibration must be validated with targets at both short and long range.
\end{sol}

\begin{sol}{prob:ch05:2}
(a)~A time offset $\delta t$ is equivalent to a displacement $v_{\mathrm{rel}}\,\delta t$: $20\times0.03=0.6$\,m for the static obstacle (relative speed $20$\,m/s) and $40\times0.03=1.2$\,m for the oncoming vehicle (relative speed $40$\,m/s). (b)~The normalised distances are $0.6/0.3=2$ and $1.2/0.3=4$, i.e. $d^{2}=4$ and $d^{2}=16$. The first is inside the gate ($4<9.21$): the detections are associated but the fused position is biased by $0.6$\,m. The second is outside ($16>9.21$): the association is rejected, each sensor spawns its own track and the oncoming vehicle appears twice, in two different positions. (c)~Augment the state with the offset $\delta t$ (a constant or slow random walk). For the delayed sensor the measurement model becomes $\mathbf{y}=\mathbf{p}(t-\delta t)+\mathbf{r}\approx\mathbf{p}(t)-\mathbf{v}_{\mathrm{rel}}\,\delta t+\mathbf{r}$, whose Jacobian with respect to $\delta t$ is $-\mathbf{v}_{\mathrm{rel}}$. The offset is observable only when the relative velocity is non-zero, and it is separable from a constant position bias only when $\mathbf{v}_{\mathrm{rel}}$ varies over time (several objects at different relative speeds, or an ego vehicle that accelerates).
\end{sol}

\begin{sol}{prob:ch05:3}
(a)~With an uncompensated bias the integrated heading is $\psi(t)=\psi_{0}+bt$, so the heading error is $b\tau$. The cross-track velocity is $v\sin(bt)\approx vbt$, and integrating gives $y(\tau)\approx vb\tau^{2}/2$. (b)~$b\tau=0.25$\,rad $=14.3^{\circ}$ and $vb\tau^{2}/2=14\times0.01\times625/2=43.8$\,m. The three-state filter of Example~\ref{ex:gnss-imu-odom} reaches $63.9$\,m and a heading error of $17.3^{\circ}$ because it enters the tunnel with a heading error that GNSS had not yet removed (its RMS error before the tunnel was already $5$\,m), because the wheel scale error adds an along-track component, and because the route turns inside the tunnel, so that the geometry is not the straight-line case of~(a). The formula gives the right order of magnitude and a lower bound. (c)~The along-track error is $s\,v\tau=0.02\times14\times25=7.0$\,m, and the total is $\sqrt{43.8^{2}+7.0^{2}}\approx44.3$\,m. A tactical-grade IMU (bias instability $0.1$--$1^{\circ}$/h, i.e. $0.5$--$5\,\mu$rad/s) reduces the cross-track term by three to four orders of magnitude but does nothing for the wheel scale error, which must be estimated as a state, as in the five-state filter, or replaced by a GNSS-derived speed.
\end{sol}

\begin{sol}{prob:ch05:4}
Script \repopath{companion\_\allowbreak repo/\allowbreak exercises/\allowbreak ch05/\allowbreak solutions/\allowbreak prob04\_\allowbreak gnss\_\allowbreak outage\_\allowbreak variants.py}. Reference run ($25$\,s outage, $2$\,Hz): error $3.56$\,m at the end of the outage, reported $3\sigma$ radius $11.94$\,m, bias $10.29\pm1.39$\,mrad/s at the tunnel entrance. (a)~$60$\,s outage: error $10.16$\,m, $3\sigma$ radius $53.62$\,m (inside), bias $10.27\pm1.39$\,mrad/s. The residual bias error of about $0.3$\,mrad/s and the residual scale error drive a drift that grows roughly with $\tau^{2}$, and the reported radius grows from $11.9$ to $53.6$\,m (a factor $4.5$ for a $2.4$ times longer outage), consistent with the quadratic law of Problem~\ref{prob:ch05:3}; the realised error ($10.2$\,m) is well inside the bound, so the filter remains consistent. (b)~$1$\,Hz fixes, $25$\,s outage: error $2.19$\,m (a single realisation; the outage error is a random variable and happens to be smaller than in the reference run), $3\sigma$ radius $13.74$\,m, bias $9.91\pm1.46$\,mrad/s and scale $3\sigma$ of $0.45$\,\% (against $0.33$\,\%). Halving the GNSS rate degrades the bias estimate only slightly because the bias is observed through the integrated effect of many fixes over the first curve, but it enlarges the scale uncertainty and the covariance predicted for the outage ($11.9\to13.7$\,m); the systematic effect of the GNSS rate is visible in the reported covariance, not in a single realised error.
\end{sol}

\begin{sol}{prob:ch05:5}
(a)~Independent rule: $\mathbf{P}_{\mathrm{f}}^{-1}=2\mathbf{P}^{-1}$, so $\mathbf{P}_{\mathrm{f}}=\mathbf{P}/2$. CI: $\mathbf{P}_{\mathrm{CI}}^{-1}=\omega\mathbf{P}^{-1}+(1-\omega)\mathbf{P}^{-1}=\mathbf{P}^{-1}$ for every $\omega$, and the estimate is $\widehat{\mathbf{x}}_{\mathrm{CI}}=\omega\widehat{\mathbf{x}}_{1}+(1-\omega)\widehat{\mathbf{x}}_{2}$. (b)~The independent rule treats the two tracks as two independent pieces of evidence and halves the variance. CI admits that they may be the same information, for example one prior relayed through two modules or two vehicles; in that case nothing has been learned by receiving it twice, and CI correctly claims no reduction of uncertainty. (c)~With $\mathbf{W}_{1}=\mathbf{W}_{2}=\tfrac{1}{2}\mathbf{I}$ the actual covariance is $\tfrac{1}{4}(\mathbf{P}+\mathbf{P}+2\rho\mathbf{P})=\tfrac{1+\rho}{2}\mathbf{P}$, so the expected NEES is $\tr\big((\mathbf{P}/2)^{-1}\tfrac{1+\rho}{2}\mathbf{P}\big)=2(1+\rho)$ for a two-dimensional state: $2$ at $\rho=0$ and $4$ at $\rho=1$. The naive rule under-reports its variance by the factor $1+\rho$, which exceeds $1.5$ for $\rho>0.5$.
\end{sol}

\begin{sol}{prob:ch05:6}
Script \repopath{companion\_\allowbreak repo/\allowbreak exercises/\allowbreak ch05/\allowbreak solutions/\allowbreak prob06\_\allowbreak ci\_\allowbreak sensitivity.py}. Over the nine combinations $\omega^{\star}$ varies only between $0.378$ (camera $5$--$30$\,\%, radar $0.5^{\circ}$) and $0.457$ (camera $5$\,\%, radar $3^{\circ}$), $\tr\mathbf{P}_{\mathrm{CI}}$ between $0.192$ and $0.211$\,m$^{2}$ and $\tr\mathbf{P}_{\mathrm{f}}$ between $0.100$ and $0.106$\,m$^{2}$, although $\tr\mathbf{P}_{\mathrm{cam}}$ ranges from $9$ to $324$\,m$^{2}$ and $\tr\mathbf{P}_{\mathrm{rad}}$ from $0.34$ to $9.9$\,m$^{2}$. The reason is that the grid only changes the \emph{weak} direction of each sensor (the camera range and the radar bearing), which carries negligible information in either rule; the fused covariance is determined by the two strong directions, the camera bearing ($0.21$\,m) and the radar range ($0.25$\,m), which are not varied. With an imaging radar of $\sigma_{\theta}=0.1^{\circ}$ the radar is better than the camera in every direction: $\omega^{\star}=0.000$, $\tr\mathbf{P}_{\mathrm{CI}}=0.073$\,m$^{2}=\tr\mathbf{P}_{\mathrm{rad}}$ (against $0.071$\,m$^{2}$ for the naive rule). When $\mathbf{P}_{2}\preceq\mathbf{P}_{1}$, every convex combination of the information matrices is dominated by $\mathbf{P}_{2}^{-1}$ alone, so CI returns the dominating track unchanged and ignores the other one; it never claims that a track that may be redundant adds information.
\end{sol}

\begin{sol}{prob:ch05:7}
(a)~$\mathrm{NDS}=\tfrac{1}{10}\big[5\times0.60+(1-0.30)+(1-0.25)+(1-0.40)+(1-0.35)+(1-0.15)\big]=\tfrac{1}{10}(3.00+3.55)=0.655$. With $\mathrm{mAVE}=0.20$ the velocity term becomes $0.80$ and $\mathrm{NDS}=\tfrac{1}{10}(3.00+3.70)=0.670$: the radar branch raises the score by $0.015$ without detecting a single additional object. (b)~$\mathrm{MOTA}=1-(120+80+10)/1000=0.79$. (c)~$\mathrm{ECE}=\tfrac{400}{1000}|0.74-0.86|+\tfrac{600}{1000}|0.91-0.95|=0.048+0.024=0.072$; the classifier is over-confident in both bins.
\end{sol}

\begin{sol}{prob:ch05:8}
(a)~$\tau_{\mathrm{e2e}}=33+5+40+5+60+10+150=303$\,ms, during which the vehicle travels $30\times0.303=9.1$\,m. (b)~The camera detection is $78$\,ms old and the radar track $15$\,ms old; both must be predicted to the same instant, the fusion time $t_{\mathrm{now}}$ (and the fused output stamped with it), i.e. forward by $78$ and $15$\,ms respectively. Over $\tau=78$\,ms the process noise of a constant-velocity model with $\sigma_{a}=0.8$\,m/s$^{2}$ adds only $\sigma_{a}\tau^{2}/2=2.4$\,mm to the position standard deviation; the dominant growth comes from the velocity uncertainty, $\sigma_{v}\tau$, e.g. $3.9$\,cm for $\sigma_{v}=0.5$\,m/s, while the predicted displacement itself, $v_{\mathrm{rel}}\tau$, is $2.3$\,m for a relative speed of $30$\,m/s and must not be omitted. (c)~Pipelining on two GPUs doubles the number of frames processed per second but leaves the latency of each result at $303$\,ms: throughput and latency are different budgets, and the reaction distance is unchanged. Shortening it requires attacking the largest terms, the actuator ($150$\,ms, e.g. brake pre-fill), the planner ($60$\,ms) and the detector ($40$\,ms, smaller network or a radar-only fast path that bypasses the camera), and the sensing interval ($33$\,ms, a $60$\,Hz camera halves it).
\end{sol}
\end{addedblock}
